\documentclass[11pt,a4paper]{article}
\usepackage{fontspec,amsmath,amssymb,amsthm,microtype}
\setmainfont{lmroman10-regular.otf}[BoldFont=lmroman10-bold.otf,ItalicFont=lmroman10-italic.otf,BoldItalicFont=lmroman10-bolditalic.otf]
\setmonofont{lmmono10-regular.otf}[Scale=MatchLowercase]
\usepackage[margin=25mm]{geometry}
\usepackage[unicode,hidelinks]{hyperref}
\setlength{\emergencystretch}{3em}
\setlength{\parskip}{0.4em}
\setlength{\parindent}{0pt}
\allowdisplaybreaks[2]
\raggedbottom
\title{\texorpdfstring{Exact velocity and clock conjugacies}{Exact velocity and clock conjugacies}}
\author{LangevinLean mathematical exposition}
\date{}
\begin{document}\maketitle
\noindent Source module: \texttt{KIPLMC/\allowbreak{}Rescaling.\allowbreak{}lean}. The accompanying HTML pairs this mathematical exposition with the exact Lean statements and records the scope of the formalization.


The physical-time particle implementation can be expressed in unit-temperature kinetic coordinates by scaling velocity and the clock. This module proves the deterministic conjugacy of the actual exponential map and the damping identity. Gaussian covariance scaling and Brownian time-change laws remain distinct probabilistic results.


\section{\texorpdfstring{Invert the velocity scale}{Invert the velocity scale}}


\texttt{rescale\allowbreak{}Velocity} sends \((v,x)\) to \((cv,x)\). If \(c\ne0\), applying the inverse scale yields \((c^{-1}cv,x)=(v,x)\). Lean simplifies the nested scalar action and cancels the nonzero factor. Position is unchanged, so a parameter-position projection commutes with this part of the change of variables.



{\raggedright\noindent Lean declarations: \texttt{rescale\allowbreak{}Velocity}, \texttt{rescale\allowbreak{}Velocity\_\allowbreak{}inverse}.\par}

\section{\texorpdfstring{Conjugate the shared exponential step}{Conjugate the shared exponential step}}


The exact coefficient scaling gives \(\psi_{c\gamma,k}(h/c)=\psi_{\gamma,k}(h)/c^k\). Apply a step with force \(g\), friction \(c\gamma\), duration \(h/c\), initial velocity \(cv\), and additive innovation whose velocity component is also scaled by \(c\). Its result equals the velocity rescaling of a physical-time step with force \(g/c^2\), friction \(\gamma\), duration \(h\). Lean proves the two scalar cancellations for the first primitive, expands the map coordinatewise, rewrites all coefficient scales, and normalizes the resulting vector identities with \texttt{module}. The innovation here is already additive; obtaining its prescribed Gaussian law under this conjugacy requires a separate covariance computation.



{\raggedright\noindent Lean declarations: \texttt{exponential\allowbreak{}Step\_\allowbreak{}conjugate}.\par}

\section{\texorpdfstring{Check the invariant friction–step product and particle scale}{Check the invariant friction–step product and particle scale}}


The OU damping depends on the product \(\gamma h\), so \(e^{-(c\gamma)(h/c)/2}=e^{-\gamma h/2}\) for \(c\ne0\). Lean unfolds the damping definition, reduces equality of exponentials to equality of exponents, and cancels the nonzero factor with \texttt{field\_\allowbreak{}simp}. For the particle choice \(c=\sqrt N\), the coefficient identity \(((\sqrt N)^2)^{-1}=N^{-1}\) follows from the nonnegative square-root square theorem. Consequently the sum-potential force in rescaled time corresponds to the average-potential force in physical time.



{\raggedright\noindent Lean declarations: \texttt{obabo\allowbreak{}Damping\_\allowbreak{}scaled}, \texttt{inverse\_\allowbreak{}sqrt\_\allowbreak{}sq}.\par}

\end{document}
